% Generated draft: requires mathematical and source review.
\documentclass[11pt]{article}
\usepackage[margin=25mm]{geometry}
\usepackage{fontspec}
\setmainfont{Latin Modern Roman}
\usepackage{amsmath,amssymb,mathtools,unicode-math}
\setmathfont{Latin Modern Math}
\usepackage{enumitem,needspace,xurl,hyperref}
\setlength{\parindent}{0pt}
\setlength{\parskip}{5pt}
\setlength{\emergencystretch}{3em}
\newcommand{\N}{\mathbb{N}}
\newcommand{\Z}{\mathbb{Z}}
\newcommand{\Q}{\mathbb{Q}}
\newcommand{\R}{\mathbb{R}}
\newcommand{\C}{\mathbb{C}}
\newcommand{\F}{\mathbb{F}}
\newcommand{\A}{\mathbb{A}}
\newcommand{\PP}{\mathbb{P}}
\newcommand{\E}{\mathbb{E}}
\newcommand{\eps}{\varepsilon}
\newcommand{\dd}{\,\mathrm{d}}
\begin{document}
\section*{UnsolvedMath ID: 30001911}
\textbf{Source problem code:} catalogue entry
\subsection{Definitions and mathematical statement}
For a nonnegative integer degree \(d\), let
\[
V_d=\bigl(\operatorname{Sym}^d(\C^3)\bigr)^*,
\]
the vector space of homogeneous degree-\(d\) polynomial forms on \(\C^3\). Its projectivization \(\PP(V_d)\) parametrizes nonzero forms up to multiplication by a nonzero scalar. A form \(f\) defines the plane curve \(\{f=0\}\subseteq\PP^2_{\C}\). The group \(G=SL_3(\C)\) acts by change of coordinates, with \((g\cdot f)(x)=f(g^{-1}x)\).

The supplied question asks whether
\[
M_d=\PP(V_d)/SL_3(\C)
\]
is rational for each of the fifteen degrees for which the source reports rationality as unknown. The record does not give those fifteen degrees. Write \(E\) for that unspecified source-dependent set, with \(|E|=15\); this notation does not identify its members or assert that precisely fifteen degrees remain unresolved today.

If the displayed quotient is intended as a rational quotient, its rational function field is
\[
K_d=\C\bigl(\PP(V_d)\bigr)^{SL_3(\C)}.
\]
Under this interpretation, the requested statement is
\[
\forall d\in E\quad \exists n_d\in\Z_{\geq0}\quad
K_d\cong_{\C}\C(t_1,\ldots,t_{n_d}),
\]
where the \(t_i\) are algebraically independent over \(\C\). Equivalently, each corresponding rational quotient is birational over \(\C\) to a projective space. Here rationality means rationality itself, rather than stable rationality.

The supplied record does not specify the quotient construction or whether a smooth, irreducible, or other restricted locus of plane curves is intended. Thus the invariant-field formulation above is an explicit conditional interpretation of the displayed quotient. Identification of \(E\) and confirmation of the intended quotient and curve locus are required to obtain a fully specified statement.
\subsection{Short English statement}
The source asks whether the moduli spaces of complex plane curves are rational in fifteen exceptional degrees. Those degrees are not listed. If the quotient means a rational quotient, the question asks whether its invariant function field is purely transcendental over \(\C\) for every degree in that missing list. The quotient convention and intended curve locus also require confirmation; the supplied open status is a dated curation claim.
\subsection{Sources}
\begin{itemize}
\item Local UnsolvedMath registry, permanent ID \(30001911\), title \emph{Rationality of Plane-Curve Moduli Spaces}, source problem code \texttt{catalogue entry}, in the Oberwolfach Reports Open Problems collection.
\item Supplied report citation: \emph{Complex Algebraic Geometry (2012)}, Oberwolfach Reports, pp.\ 2669--2733, DOI \texttt{10.4171/owr/2011/47}. \url{https://doi.org/10.4171/owr/2011/47}. The record gives proposed year \(2011\) and does not identify a proposer.
\item Supplied literature source: ``Bohning, rationality problem in invariant theory.'' \url{https://warwick.ac.uk/fac/sci/maths/people/staff/boehning/invariant_theory_book.pdf}.
\item Supplied literature source: ``original OWR report.'' \url{https://doi.org/10.4171/owr/2011/47}.
\item Literature checked at: \texttt{2026-08-21}. The supplied dated curation labels the problem open and reports a finite exceptional list of degrees with unsettled rationality, so that the extracted question is not closed in full generality. This assessment is supplied dated curation, not independently verified here; the record itself requests independent verification of time-sensitive status and cited claims.
\item The supplied record labels the statement \texttt{corrected\_verified} and says that its mathematical content was checked against the official Oberwolfach report, with notation, context, and question form normalized without changing the source-level claim. This is a supplied verification claim, not an independent verification in this draft.
\end{itemize}
\end{document}
